% Ad Familiares 7.20 (ad-familiares-07-20)
% Source: https://raw.githubusercontent.com/PerseusDL/canonical-latinLit/master/data/phi0474/phi056/phi0474.phi056.perseus-lat1.xml
% Fetched by scripts/fetch_latin.py

% Dateline (Perseus): Scr. Veliae xiii K. Sext. a. 710 (44).

\ciceroLetterOpener{CICERO TREBATIO S.}

\ciceroSection{1} amabilior mihi Velia fuit, quod te ab ea sensi amari. sed quid ego dicam 'te,' quem quis non amat? rufio medius fidius tuus ita desiderabatur ut si esset unus e nobis. sed te ego non reprehendo qui illum ad aedificationem tuam traduxeris. quamquam enim Velia non est vilior quam Lupercal, tamen istuc malo quam haec omnia. tu si me audies, quem soles, has paternas possessiones tenebis (nescio quid enim Velienses verebantur) neque Haletem, nobilem amnem, relinques nec Papirianam domum deseres; quamquam illa quidem habet lotum, a quo etiam advenae teneri solent; quem tamen si excideris, multum prospexeris.

\ciceroSection{2} sed in primis opportunum videtur, his praesertim temporibus, habere perfugium primum eorum urbem quibus carus sis, deinde tuam domum tuosque agros, eaque remoto, salubri, amoeno loco; idque etiam mea interesse, mi Trebati, arbitror. sed valebis meaque negotia videbis meque dis iuvantibus ante brumam exspectabis.

\ciceroSection{3} ego a Sex. Fadio, Niconis discipulo, librum abstuli *ni/kwnos peri\ *polufagi/as . O medicum suavem meque docilem ad hanc disciplinam! sed Bassus noster me de hoc libro celavit, te quidem non videtur. ventus increbrescit. cura ut valeas. xiii K. Sext. Velia.
